\scanpage{49}
\mspage{136}

$\underline{\mathrm{Ch}}_{\infty}(S,s_0,-)$ par le foncteur but
\[
\unsure{\beta|b}:\underline{\mathrm{Ch}}_{\infty}(S)\longrightarrow S\quad).
\]

On a prouvé \deleted{\unsure{dans}} loc. cit. que $\bar f$ est une
Hot-fibration (``propre'' et ``lisse''), mais dans un
contexte qui semblait faire appel tacitement à
la théorie du dérivateur $\mathrm{HOT}$. Donc pour
mettre les résultats ``topologiques'' de loc. cit.
sur une base solide, il faut développer les
notions pertinentes de foncteur propre, foncteur
lisse, sans référence à une théorie $\mathrm{HOT}$. C'est
ce que j'avais d'ailleurs fait déjà \deleted{\unsure{dans}}
en 1983, du temps de \unsure{Pursuing} Stacks,
mais \deleted{sans} rédaction en forme à ce sujet. Il
me faut donc à présent revenir sur ces
notions, de façon un peu systématique. D'ailleurs,
j'aurais du mal à construire un seul foncteur
\deleted{\unsure{propre et lisse}}\inserted{Hot-fibrant} (\unsure{s.f.|sauf} \deleted{\unsure{en}} les exemples triviaux,
tels les isomorphismes\deleted{\unsure{)}}, ou les produits $Y\times Z\xrightarrow{\mathrm{pr}_1}Y$),
en disposant d'une théorie des
foncteurs lisses et des foncteurs propres.$^{*}$
% slanted
\margin{$^{*}$ J'exagère un peu, quand même -- après tout il doit être \unsure{assez évident|clair} que les foncteurs \unsure{qui sont des fibrés au sens de|ou les fibrés au sens de|fibrés au sens de} Kan sont des \deleted{\unsure{Hot-}}\unsure{Hot-fibrés|Hot-fibrants|Hot-fibrations} \unsure{pour le faire|sans faire|pas des tas} (mais il y en a quand même peu!).}

% archive sheet number: 49
